\section{The conforming Galerkin method}
\label{sec:galerkin}

Let $\{\Th\}_{h>0}$ be a shape-regular family of conforming triangulations of
$\Omega$, and let $h=\max_{K\in\Th}\operatorname{diam}(K)$. For continuous,
piecewise affine elements define
\begin{equation}
  \Vh
  :=\left\{v_h\in C^0(\overline\Omega):
  v_h|_K\in\mathbb P_1(K)\ \forall K\in\Th,
  \quad v_h|_{\partial\Omega}=0\right\}.
  \label{eq:finite-element-space}
\end{equation}
Because $\Vh\subset\V$, the approximation is conforming.

The discrete problem is to find $u_h\in\Vh$ such that
\begin{equation}
  a(u_h,v_h)=F(v_h)
  \qquad\text{for every }v_h\in\Vh.
  \label{eq:discrete-problem}
\end{equation}
The same continuity and coercivity constants as in the continuous problem
apply on $\Vh$, so existence and uniqueness follow immediately.

\begin{proposition}[Galerkin orthogonality]
\label{prop:galerkin-orthogonality}
Let $u$ and $u_h$ solve \cref{eq:weak-problem,eq:discrete-problem}. Then
\begin{equation}
  a(u-u_h,v_h)=0
  \qquad\text{for every }v_h\in\Vh.
  \label{eq:galerkin-orthogonality}
\end{equation}
\end{proposition}

\begin{proof}
Restrict \cref{eq:weak-problem} to test functions in $\Vh$ and subtract
\cref{eq:discrete-problem}.
\end{proof}

If $\{\phi_i\}_{i=1}^N$ is the nodal basis of $\Vh$ and
$u_h=\sum_{j=1}^N U_j\phi_j$, then \cref{eq:discrete-problem} becomes
\begin{equation}
  \bm K\bm U=\bm b,
  \qquad
  K_{ij}=a(\phi_j,\phi_i),
  \quad
  b_i=F(\phi_i).
  \label{eq:linear-system}
\end{equation}
The matrix $\bm K$ is symmetric positive definite, permitting a conjugate
gradient solve with an appropriate preconditioner.

\begin{remark}[Energy projection]
The consistency identity in \cref{eq:galerkin-orthogonality} shows that $u_h$
is the projection of $u$ onto $\Vh$ in the energy inner product. This geometric
interpretation is the starting point of the quasi-optimality estimate.
\end{remark}

% DEMO: Insert a displayed formula here, for example an energy projection identity.
